\begin{frame}[fragile]{The Bottleneck}
  \begin{columns}[T]
    \column{0.38\textwidth}
      \textbf{\texttt{count\_up\_down-1}}\\[0.4em]
      \small
\begin{verbatim}
int x = 0, y = n;
while (x < n) {
    x++;  y--;
}
assert(y == 0);
\end{verbatim}
      \vspace{0.2em}
      \begin{tabular}{ll}
        \toprule
        Refinements & \alert{69} \\
        Wall time   & \alert{75\,s} \\
        Result      & \alert{UNKNOWN} \\
        \bottomrule
      \end{tabular}

    \column{0.58\textwidth}
      \textbf{CEGAR finds (per round):}\\[0.4em]
      {\small\ttfamily
      \textcolor{WarnRed}{y=0,\ x=n}\\[0.15em]
      \textcolor{WarnRed}{y=1,\ x=n{-}1}\\[0.15em]
      \textcolor{WarnRed}{y=2,\ x=n{-}2}\\[0.15em]
      \hspace*{0.5em}\(\vdots\)\quad one per unrolling}\\[0.7em]

      \textbf{LLM guesses:}\\[0.4em]
      {\small\ttfamily
      \textcolor{GoodGreen}{x + y = n}\quad\small\textcolor{GoodGreen}{(loop invariant)}\\[0.15em]
      \textcolor{GoodGreen}{x \(\geq\) 0},\ \textcolor{GoodGreen}{y \(\leq\) n}}
  \end{columns}
  \note{The canonical CEGAR failure on loops. x counts up, y counts down — the invariant x+y=n holds every iteration but interpolation never sees it because each counterexample only reveals one step. CEGAR finds instance-specific predicates: y=0,x=n at step 0; y=1,x=n-1 at step 1; etc. — one per unrolling, never the general invariant. LLM proposes x+y=n directly from the source code. After injection CEGAR converges in 2 steps.}
\end{frame}
